
\chapter{Станционная постановка: данные, признаки и базовые модели}

\section{Общая схема экспериментального исследования}

В рамках данной работы экспериментальное исследование строится как последовательное решение задачи прогнозирования метеорологических характеристик по пространственно-временным данным. Исходной основой выступает набор RusWeather-GF, содержащий погодные поля и связанные с ними наблюдения, представленные в форме, пригодной для машинного обучения. Основная цель экспериментов состоит в том, чтобы по историческому фрагменту климатических данных восстановить закономерности их изменения и получить прогноз целевых переменных на заданном горизонте.

Общая схема исследования включает несколько логически связанных этапов. На первом этапе выполняется подготовка данных: выбираются используемые метеорологические признаки, формируются пространственно-временные последовательности наблюдений, данные приводятся к согласованному формату и разделяются на обучающую, валидационную и тестовую выборки. Такой этап необходим для того, чтобы обеспечить корректное обучение моделей и сопоставимость результатов.

На втором этапе формулируется задача прогноза как задача отображения истории наблюдений в будущее состояние системы. Иными словами, модель получает на вход некоторый временной интервал предшествующих измерений и должна предсказать значения интересующих погодных характеристик на следующих шагах. В зависимости от постановки эксперимента прогноз может рассматриваться как детерминированный или вероятностный, что особенно важно при анализе экстремальных погодных явлений.

На третьем этапе обучаются глубокие нейросетевые модели, способные учитывать как временные зависимости в данных, так и пространственную структуру метеорологических полей. Сравнение моделей проводится в единой экспериментальной постановке, что позволяет оценивать вклад архитектурных решений не изолированно, а в контексте общей задачи прогноза.

На заключительном этапе выполняется оценка качества полученных результатов. Рассматриваются стандартные метрики точности прогноза, а также характеристики, связанные с устойчивостью модели к редким и аномальным наблюдениям. Отдельное внимание уделяется анализу экстремальных ситуаций, поскольку именно в таких режимах наиболее важны практическая применимость модели и надёжность её выводов.

\section{Станционные данные RusWeather-GF}

В станционной части работы используется набор данных RusWeather-GF. Каждая запись содержит идентификатор станции, дату наблюдения, значения температуры и осадков, а также географические характеристики станции: широту, долготу и высоту. В качестве целевой переменной используется \texttt{precipitation}, а температура рассматривается как дополнительный динамический предиктор.

Частота наблюдений определяется автоматически как 24 часа. Это позволяет естественным образом формулировать целевые горизонты:
\[
\text{TP24},\ \text{TP72},\ \text{TP120},\ \text{TP240},
\]
что соответствует прогнозу суммарных осадков на \(1\), \(3\), \(5\) и \(10\) суток вперед. Такая постановка удобна по двум причинам. Во-первых, горизонты охватывают как краткосрочный, так и среднесрочный режим прогноза. Во-вторых, они хорошо согласуются с практикой предупреждения опасных осадочных событий, для которых важны как ближайшие сутки, так и несколько суток вперед.

Объем исходного станционного массива достаточно велик: порядка 8.9 млн строк наблюдений. После временного разбиения с учетом purge-горизонта получены следующие объемы:
\begin{center}
\begin{tabular}{lrr}
Разбиение & Число строк & Комментарий \\
Train & 6\,409\,702 & обучающая часть \\
Validation & 1\,276\,786 & валидационная часть \\
Test & 1\,196\,665 & тестовая часть \\
\end{tabular}
\end{center}

\begin{table}[H]
\centering
\begin{tabular}{p{0.38\textwidth}p{0.50\textwidth}}
Параметр & Значение \\
Целевая переменная & сумма осадков \\
Дополнительный динамический признак & температура \\
Горизонты прогноза & 24, 72, 120, 240 часов \\
Лаги & 1, 3, 6, 12, 24, 48, 72, 120, 240 часов \\
Скользящие окна & 3, 6, 12, 24, 72, 120, 240 часов \\
Порог экстремальности & \(q_{0.9}\) \\
Минимум наблюдений для порога & 50 на станцию \\
Доли train/val/test & 0.70 / 0.15 / 0.15 \\
Защитный интервал & 240 часов \\
\end{tabular}
\caption{Основные параметры подготовки станционных данных}
\label{tab:cfg-station-main}
\end{table}

Использование purge-интервала в 240 часов позволяет избежать утечки информации между соседними участками времени при многогоризонтном прогнозе. Иначе наблюдения, близкие к границе train/validation или validation/test, могли бы использовать близкие целевые значения и тем самым искусственно завышать качество.

\section{Построение признаков для станционного прогноза}

Одним из центральных элементов станционной постановки является инженерия признаков. В работе используется несколько классов признаков.

\subsection{Лаговые признаки}

Для осадков и температуры формируются лаги с набором горизонтов
\[
1,\ 3,\ 6,\ 12,\ 24,\ 48,\ 72,\ 120,\ 240 \text{ часов}.
\]
Поскольку фактический шаг данных составляет 24 часа, часть этих значений сопоставляется доступной сетке времени. Лаговые признаки отражают краткосрочную инерцию и сезонно-синоптическую структуру процесса. Для осадков это особенно важно: наличие осадков в предшествующие дни повышает вероятность сохранения влажного режима, тогда как после сухого периода вероятность экстремума ниже.

\subsection{Скользящие статистики}

Для тех же окон времени строятся скользящие статистики: среднее, сумма и стандартное отклонение. Сумма по окну интерпретируется как интегральная характеристика накопленной влаги, среднее — как усредненный фон, а стандартное отклонение — как индикатор нестабильности процесса. В задачах экстремальных осадков такие статистики особенно полезны, поскольку экстремум часто возникает не изолированно, а на фоне уже аномального режима.

\subsection{Календарные признаки}

Используются признаки месяца, а также day-of-year. Это позволяет модели учитывать сезонность. Для осадков сезонная компонента критична: характер распределения и интенсивность осадков существенно различаются между холодным и теплым сезоном, а также по регионам. Календарные признаки дают модели компактное представление этой сезонной структуры.

\subsection{Географические признаки}

В качестве статических признаков используются широта, долгота и высота станции. Эти переменные частично компенсируют отсутствие явной пространственной структуры в табличной постановке. Например, горные районы и приморские области обладают иным режимом осадков, чем континентальные равнинные территории. В станционной модели координаты не создают полноценной карты, но дают алгоритму возможность различать регионы с разной климатической нормой.

\section{Целевые переменные и событийное преобразование}

Для каждого временного шага формируются целевые переменные \(\mathrm{TP24}\), \(\mathrm{TP72}\), \(\mathrm{TP120}\), \(\mathrm{TP240}\), представляющие будущие суммы осадков на соответствующих горизонтах. Модели обучаются в регрессионной постановке, то есть минимизируют ошибку по непрерывной величине.

Однако сама логика работы требует анализа не только непрерывного прогноза, но и способности модели выделять аномальные случаи. Поэтому валидация строится через высококвантильные события. Для каждой станции и каждого горизонта оценивается порог \(q_{0.9}\) при условии достаточного числа наблюдений. Далее для фактического значения формируется бинарная метка превышения порога, а из непрерывного прогноза выводится вероятность события через температурно-сглаженное сравнение с \(q_{0.9}\).

Такая схема имеет несколько достоинств. Во-первых, она позволяет использовать одну и ту же регрессионную модель одновременно для прогноза величины и для детекции аномалий. Во-вторых, она локально адаптивна: порог экстремальности зависит от станции. В-третьих, метрики ROC-AUC, ROCSS, PR-AUC и Brier дают более содержательную картину для прикладного использования, чем усредненная квадратичная ошибка.

\section{Baseline-модели}

\subsection{Климатология}

Климатологическая базовая модель строится как среднее значение осадков по станции и месяцу. Для каждого горизонта вычисляется станционно-месячное среднее, которое и используется как прогноз. Несмотря на простоту, климатология является сильным ориентиром, потому что она уже учитывает устойчивую сезонную компоненту и региональную специфику станции.

\subsection{SGD-регрессия}

В качестве линейной базовой модели использована модель \texttt{SGDRegressor} с квадратичной функцией потерь и elastic net-регуляризацией. Обучение организовано по схеме \texttt{partial\_fit} с несколькими эпохами прохода по шардированным данным. Эта модель важна как интерпретируемая отправная точка: если даже линейная модель показывает удовлетворительное качество, сложные алгоритмы должны оправдывать свою стоимость приростом метрик.

\subsection{Градиентные бустинги}

Основными табличными моделями являются LightGBM и XGBoost. Обе модели обучаются инкрементно блоками по 100 итераций до общего лимита 500 раундов. Такой режим согласуется с выбранной метрикой отбора \(\mathrm{mean\_rocss\_val}\), то есть средним ROCSS по горизонтам валидации. Использование именно ROCSS как селекционной метрики принципиально: оно задает приоритет способности модели детектировать экстремальные события, а не только минимизировать регрессионную ошибку.

\begin{table}[H]
\centering
\begin{tabular}{p{0.30\textwidth}p{0.28\textwidth}p{0.28\textwidth}}
Параметр & LightGBM & XGBoost \\
Функция потерь & регрессия & квадратичная ошибка \\
Скорость обучения & 0.05 & 0.05 \\
Число итераций & 500 & 500 \\
Глубина / листья & 31 лист & глубина 6 \\
Подвыборка объектов & 0.9 & 0.9 \\
Подвыборка признаков & 0.9 & 0.9 \\
Случайное зерно & 42 & 42 \\
\end{tabular}
\caption{Основные параметры бустинговых моделей}
\label{tab:boosting-params-main}
\end{table}

\section{Результаты станционного эксперимента}

В таблице~\ref{tab:station-mean} приведены усредненные событийные метрики на валидационной выборке для всех реализованных станционных моделей. Численные значения восстановлены непосредственно из результатов ноутбука.

\begin{table}[H]
\centering
\begin{tabular}{lcccc}
Модель & mean ROC-AUC & mean ROCSS & mean PR-AUC & mean Brier \\
Климатология & 0.7062 & 0.4124 & 0.2023 & 0.1232 \\
SGD & 0.4666 & -0.0667 & 0.0927 & 0.1658 \\
LightGBM & 0.7213 & 0.4425 & 0.2158 & 0.1249 \\
XGBoost & \textbf{0.7235} & \textbf{0.4469} & \textbf{0.2170} & 0.1247 \\
\end{tabular}
\caption{Средние событийные метрики на валидации для станционных моделей}
\label{tab:station-mean}
\end{table}

Из таблицы видны несколько принципиальных выводов.

Во-первых, климатология оказывается сильной базовой моделью. Ее средний ROCSS равен \(0.4124\), что существенно выше случайного уровня и указывает на высокую значимость сезонно-климатической компоненты в данных. Это важный результат: любая сложная модель в задачах осадков должна превосходить не \enquote{нулевой} базовая модель, а именно сильную климатологию.

Во-вторых, линейная модель SGD демонстрирует неудовлетворительное качество. Средний ROCSS отрицателен, а Brier score заметно хуже, чем у остальных моделей. Это говорит о том, что зависимость между признаками и целевыми осадками не может быть адекватно описана простой линейной комбинацией признаков. Даже богатая инженерия признаков не компенсирует отсутствие нелинейности.

В-третьих, ансамблевые бустинги уверенно превосходят климатологию и линейную модель. Между LightGBM и XGBoost разрыв невелик, однако XGBoost дает лучшие средние значения ROC-AUC, ROCSS и PR-AUC. Следовательно, в рамках станционной постановки именно градиентные бустинги представляют наилучший компромисс между качеством и вычислительной реализуемостью.

Для более детального анализа в таблице~\ref{tab:station-horizon} приведены результаты по отдельным горизонтам.

\begin{table}[H]
\centering
\begin{tabular}{llcccc}
Горизонт & Модель & ROC-AUC & ROCSS & PR-AUC & Brier \\
\multirow{4}{*}{24 ч}
& Климатология & 0.6267 & 0.2535 & 0.1512 & 0.1394 \\
& SGD & 0.5140 & 0.0280 & 0.1067 & 0.1625 \\
& LightGBM & 0.6871 & 0.3743 & 0.1986 & 0.1374 \\
& XGBoost & \textbf{0.6905} & \textbf{0.3809} & \textbf{0.2007} & \textbf{0.1370} \\
\multirow{4}{*}{72 ч}
& Климатология & 0.7022 & 0.4044 & 0.1962 & 0.1184 \\
& SGD & 0.4638 & -0.0724 & 0.0911 & 0.1588 \\
& LightGBM & 0.7131 & 0.4262 & 0.2084 & 0.1202 \\
& XGBoost & \textbf{0.7151} & \textbf{0.4302} & \textbf{0.2093} & \textbf{0.1200} \\
\multirow{4}{*}{120 ч}
& Климатология & 0.7294 & 0.4587 & 0.2161 & \textbf{0.1167} \\
& SGD & 0.4514 & -0.0972 & 0.0880 & 0.1650 \\
& LightGBM & 0.7285 & 0.4570 & 0.2186 & 0.1195 \\
& XGBoost & \textbf{0.7298} & \textbf{0.4596} & \textbf{0.2192} & 0.1194 \\
\multirow{4}{*}{240 ч}
& Климатология & \textbf{0.7664} & \textbf{0.5329} & \textbf{0.2458} & \textbf{0.1184} \\
& SGD & 0.4373 & -0.1253 & 0.0848 & 0.1766 \\
& LightGBM & 0.7563 & 0.5126 & 0.2375 & 0.1224 \\
& XGBoost & 0.7584 & 0.5169 & 0.2385 & 0.1222 \\
\end{tabular}
\caption{Валидационные метрики по отдельным горизонтам прогноза}
\label{tab:station-horizon}
\end{table}

Эта таблица позволяет сделать более тонкие наблюдения.

На горизонте 24 часа бустинги дают явное преимущество над климатологией. Здесь свежая динамика и нелинейное сочетание признаков наиболее информативны, поэтому XGBoost и LightGBM существенно улучшают ранжирование экстремальных случаев.

На горизонтах 72 и 120 часов преимущество бустингов сохраняется, но становится умеренным. Можно сделать вывод, что модель уже больше опирается на медленные режимные закономерности, и вклад сложной нелинейности снижается. Тем не менее XGBoost по-прежнему удерживает наилучшие значения ROC-AUC и ROCSS.

На горизонте 240 часов картина меняется: климатология становится сильнейшей моделью по всем основным событийным метрикам. Это очень важный эмпирический вывод. Для дальнего горизонта станционные признаки из прошлого уже не обеспечивают достаточно устойчивой информации, и доминировать начинает долгосрочная климатическая структура по станции и сезону. Следовательно, глубокие и сложные модели не обязаны превосходить климатологию на дальних горизонтах, если в них отсутствует явное пространственное представление крупномасштабной циркуляции.

\section{Интерпретация результатов станционного блока}

Полученные результаты позволяют сформулировать несколько выводов, значимых не только для данной работы, но и для общей постановки прогнозирования осадков.

Первый вывод состоит в том, что инженерия признаков действительно работает. Даже без нейросетевых методов табличное представление на станционном уровне способно обеспечить хорошее качество событийного ранжирования. Это означает, что для прикладных систем локального прогноза бустинговые модели представляют собой разумную и сильную базовую модель.

Второй вывод связан с ограничениями станционного подхода. Хотя бустинги хорошо работают на коротких и средних горизонтах, они не моделируют пространственную структуру осадков в явном виде. Любая зависимость между соседними районами выражается только косвенно через координаты, а это ограничивает способность модели улавливать крупномасштабные атмосферные процессы. Именно поэтому для сеточного прогноза и пространственной детекции аномалий требуется другой класс моделей.

Третий вывод касается самого понятия аномалии. В станционной части работы аномалия успешно формализуется как локальное превышение квантильного порога. Это позволяет перейти от регрессионной задачи к вероятностной без смены модели и без необходимости строить отдельный классификатор. Такая формализация будет использована и в сеточной части, но уже на уровне пространственных карт \(q_{0.9}\).

\section{Роль станционной базовой модели в общей структуре работы}

Несмотря на то что официальная тема работы фокусируется на глубоких нейросетевых моделях, станционная глава имеет принципиальное значение. Она выполняет три функции.

Во-первых, она задает нижнюю и среднюю планку качества. Сильный бустинговый базовая модель необходим для того, чтобы последующее использование глубокой модели не выглядело неоправданным усложнением.

Во-вторых, она позволяет отделить эффект \enquote{более богатых данных} от эффекта \enquote{более сложной архитектуры}. Если одна лишь инженерия признаков и нелинейный табличный алгоритм уже дают хороший результат, значит, дальнейшее улучшение должно происходить за счет перехода к пространственно-согласованному представлению данных.

В-третьих, станционный блок задает язык интерпретации. Через него удобно объяснить, почему событийные метрики, квантильные пороги и вероятностная оценка редких событий более релевантны прикладной задаче, чем только RMSE или MAE. В дальнейшем та же логика переносится на сеточный прогноз, где аномалии уже детектируются как экстремальные осадки по картам квантилей.

Таким образом, станционная базовая модель не является в работе побочным сюжетом. Она выступает как методологическая основа, на фоне которой можно содержательно оценивать вклад глубоких нейросетевых моделей в задачах прогнозирования и детекции аномалий в климатических данных.
